\subsection{Abelian categories and exact sequences}
\label{sec:p10-abelian-exact}

Freyd's structural definition starts with universal constructions and
normality, rather than assuming that arrows can already be added.
The development below follows that order: images precede biproducts,
biproducts provide a commutative addition, and an invertible shear supplies
additive inverses. The locators refer to the official TAC reprint of
\emph{Abelian Categories}, whose original publication was in 1964
\cite{Freyd2003TAC}. The hom-group definition of an additive category is
then explicitly connected to this structural presentation.

We use the locally small categories, composition convention, and size
framework of Definition~\ref{def:p10-category} and
Convention~\ref{def:p10-categorical-size}. No smallness of the entire
category, faithful module embedding, or supply of projectives/injectives
is assumed. A subobject is an equivalence class of monomorphisms into
a fixed object; equality of subobjects means a compatible isomorphism,
not equality of chosen carrier objects. Quotient objects are the dual
classes of epimorphisms out of a fixed object.

\begin{hbdefinition}{def:p10-ab-zero-kernel}{Zero objects, kernels, and cokernels}
A zero object \(0\) is both initial and terminal. The composite
\(A\to0\to B\) is the zero arrow \(0_{A,B}\); its composites with
any compatible arrows are again zero arrows. A kernel of \(f:A\to B\)
is an arrow \(k:K\to A\) with \(fk=0\), universal among arrows
\(x:T\to A\) with \(fx=0\). A cokernel is an arrow \(q:B\to Q\)
with \(qf=0\), universal among arrows out of \(B\) killing \(f\).
Write \(\operatorname{Ker}f\) and \(\operatorname{Cok}f\) for these
chosen arrows or their source/target objects, with the type stated.
They are the equalizer/coequalizer of \(f\) and the zero arrow, not
the kernel/cokernel of a set function without its categorical structure.
\end{hbdefinition}

\begin{hbdefinition}{def:p10-ab-freyd}{Freyd's structural abelian axioms}
A locally small category is abelian here if it satisfies the source's
axioms on printed p.~35: A0, a zero object; A1 and A1*, binary products
and coproducts; A2 and A2*, kernels and cokernels of every arrow; A3,
every monomorphism is a kernel of some arrow; A3*, every epimorphism
is a cokernel of some arrow. A3 and A3* are normality conditions.
They do not assert that a specific arrow's image/coimage comparison is
invertible; that conclusion is proved below. The opposite category
satisfies the same list, with the starred/unstarred axioms interchanged.
\end{hbdefinition}

\begin{hbtheorem}{thm:p10-ab-normal}{proved=annex:p10.ab.1}{Normality, intersections, and equalizers}
Every monomorphism is the kernel of its own cokernel, and every
epimorphism is the cokernel of its own kernel. An arrow both monic and
epic is invertible. Two subobjects have a pullback intersection;
every pair of parallel arrows has an equalizer. Dually there are
coequalizers, pushouts, and binary subobject unions. Hence all finite limits and colimits
exist, by Theorem~\ref{thm:p10-limits-products-equalizers}.
This expands the universal-construction arguments of source
Theorems~2.11--2.15, pp.~36--42, before using addition.
\end{hbtheorem}

\begin{hbdefinition}{def:p10-ab-images}{Image and coimage}
The image of \(f:A\to B\) is the least subobject of \(B\) allowing
a factorization of \(f\). The coimage is the dual least quotient
through which \(f\) factors. The candidates are
\[
 \operatorname{Im}f=\operatorname{Ker}(\operatorname{Cok}f),\qquad
 \operatorname{Coim}f=\operatorname{Cok}(\operatorname{Ker}f).
\]
The arrows \(q_f:A\to\operatorname{Coim}f\) and
\(i_f:\operatorname{Im}f\to B\) are part of this data.
\end{hbdefinition}

\begin{hbtheorem}{thm:p10-ab-factorization}{proved=annex:p10.ab.2}{The canonical comparison is an isomorphism}
These candidates have the stated universal properties. The factorization
\[
 A\xrightarrow{q_f}\operatorname{Coim}f
 \xrightarrow{\bar f}\operatorname{Im}f\xrightarrow{i_f}B
\]
has \(q_f\) epic, \(i_f\) monic, and \(\bar f\) an isomorphism.
Every epi--mono factorization has a unique compatible isomorphism to
this one. Moreover \(f\) is monic iff its kernel is zero, and epic iff
its cokernel is zero. These are source Theorems~2.16--2.19,
pp.~42--44; no module elements are used in their proofs.
\end{hbtheorem}

\begin{hbtheorem}{thm:p10-ab-biproduct}{proved=annex:p10.ab.3}{Constructing the biproduct, before addition}
The canonical map \(A+B\to A\times B\), with identity diagonal
entries and zero off-diagonal entries, is an isomorphism. Write the
resulting object \(A\oplus B\), with injections \(u_1,u_2\) and
projections \(p_1,p_2\). It has both universal properties and
\(p_i u_j=1\) for \(i=j\), \(0\) otherwise. This is the source's
direct-sum construction, Theorems~2.31--2.35, pp.~45--47.
\end{hbtheorem}

For matrices we retain Freyd's row-source convention: the entry in row
\(i\), column \(j\) of a map \(\bigoplus_i A_i\to\bigoplus_j B_j\)
is an arrow \(A_i\to B_j\), namely \(p_jfu_i\).

\begin{hbtheorem}{thm:p10-ab-addition}{proved=annex:p10.ab.4}{Addition and negatives are derived, not postulated}
Let \(\Delta_A:A\to A\oplus A\) be the product diagonal and
\(\nabla_B:B\oplus B\to B\) the coproduct codiagonal. Set
\[
 f+g=\nabla_B(f\oplus g)\Delta_A\qquad(f,g:A\to B).
\]
Each \(\mathcal C(A,B)\), also denoted \((A,B)\) in Freyd's
notation, is an abelian group with zero the zero arrow. Composition
is bilinear, and the operation is uniquely determined by the categorical
biproducts. In particular the shear
\(\left(\begin{smallmatrix}1&f\\0&1\end{smallmatrix}\right)\)
is proved invertible without already assuming a negative of \(f\);
its inverse supplies that negative. This reconstructs
Propositions~2.36--2.37 and Theorems~2.38--2.39, pp.~47--49.
\end{hbtheorem}

\begin{hbdefinition}{def:p10-ab-additive}{Preadditive and additive categories}
A preadditive category has an abelian group structure on each hom set
and bilinear composition. In Freyd's Exercise~2-A, p.~60, a zero object
is also included in the convention. An additive category in that
convention has finite products, equivalently finite biproducts.
We retain that zero-object convention here. In a given preadditive
category, addition is structure; in an abelian category it is the
canonically derived structure of the preceding theorem.
\end{hbdefinition}

\begin{hbtheorem}{thm:p10-ab-additive-bridge}{proved=annex:p10.ab.5}{The additive presentation and biproduct equations}
In a preadditive category with zero object and binary products, those
products are also coproducts. A direct-sum system is equivalently
specified by
\[
 p_i u_j=\delta_{ij},\qquad u_1p_1+u_2p_2=1.
\]
Thus the additive definition with kernels, cokernels, and normal monos
and epis is equivalent to the structural abelian definition above.
The prescribed hom addition agrees with the canonical addition.
Matrices of arrows between finite biproducts compose by the ordinary
sum-of-composites formula. The statements include empty biproducts.
\end{hbtheorem}

\begin{hbdefinition}{def:p10-ab-exact}{Exact and short exact sequences}
A composable pair \(A\xrightarrow f B\xrightarrow g C\) is exact
at \(B\) when \(\operatorname{Im}f=\operatorname{Ker}g\) as
subobjects of \(B\). A sequence is exact when this holds at each
displayed interior object. A short exact sequence is
\(0\to A\xrightarrow i B\xrightarrow p C\to0\), exact at all
three nonzero positions; the symbols \(A,B,C\) need not be nonzero objects.
\end{hbdefinition}

\begin{hbtheorem}{thm:p10-ab-exactness}{proved=annex:p10.ab.6}{Exactness criteria and the endpoint cases}
For a pair \((f,g)\), exactness is equivalent to
\(gf=0\) and \(q_f k_g=0\), where \(k_g\) is a kernel of \(g\)
and \(q_f\) a cokernel of \(f\). It is also equivalent to
\(\operatorname{Coim}g=\operatorname{Cok}f\) as quotients of \(B\).
A short sequence is exact iff \(i=\operatorname{Ker}p\) and
\(p=\operatorname{Cok}i\), up to their unique compatible isomorphisms.
The sequence \(0\to A\xrightarrow f B\) is exact iff \(f\) is
monic; \(A\xrightarrow f B\to0\) is exact iff \(f\) is epic;
\(0\to A\xrightarrow f B\to0\) is exact iff \(f\) is invertible.
Every arrow has the short exact kernel--image and image--cokernel
sequences coming from its factorization.
\end{hbtheorem}

\begin{hbtheorem}{thm:p10-ab-splitting}{proved=annex:p10.ab.7}{Splitting a short exact sequence}
For \(0\to A\xrightarrow i B\xrightarrow p C\to0\), these are
equivalent: \(p\) has a section; \(i\) has a retraction; there is a
compatible isomorphism \(B\cong A\oplus C\).
For a section \(s\), the retraction \(r\) is uniquely specified by
\(ir=1-sp\); it satisfies \(ri=1\) and \(rs=0\).
The maps \((i,s)\) and \((r,p)\) are inverse biproduct maps.
This gives an explicit proof of the source's splitting statement 2.68,
p.~59, without treating every short exact sequence as split.
\end{hbtheorem}

\begin{hbtheorem}{thm:p10-ab-base-change}{proved=annex:p10.ab.8}{Pullbacks, pushouts, and stable exactness}
The pullback of \(f:A\to D\), \(g:B\to D\) is the kernel of
\((f,-g):A\oplus B\to D\), with the induced projections.
Dually a pushout is a cokernel of the signed injection pair.
Pullbacks preserve epimorphisms and their kernel sequences; dually,
pushouts preserve monomorphisms and their cokernel sequences.
A square \(P\to B\), \(P\to A\), \(A\to D\), \(B\to D\)
is a pullback iff its signed sequence
\(0\to P\to A\oplus B\to D\) is exact, and a pushout iff
\(P\to A\oplus B\to D\to0\) is exact. Both conditions hold
precisely when that sequence is short exact.
\end{hbtheorem}

\begin{hbtheorem}{thm:p10-ab-local-arrows}{proved=annex:p10.ab.9}{A legitimate generalized-arrow method}
For \(f:A\to B\) and \(x:T\to B\), \(x\) factors through
\(\operatorname{Im}f\) iff there are an epimorphism \(s:E\to T\)
and an arrow \(y:E\to A\) with \(fy=xs\).
Equality of two arrows can be tested after any epic precomposition;
successive epic covers compose. Thus kernel/image containment may be
proved by such local lifts, including a lift of the universal kernel arrow
itself. This method is derived from pullbacks and normality, not from
an unproved embedding into modules or an assumed underlying set of elements.
\end{hbtheorem}

\begin{hbtheorem}{thm:p10-ab-short-five}{proved=annex:p10.ab.10}{The short five lemma}
In a commutative diagram of short exact sequences, with vertical maps
\(a:A\to A'\), \(b:B\to B'\), \(c:C\to C'\), if \(a\) and
\(c\) are isomorphisms, so is \(b\). The proof separately establishes
monicity and epicity by kernel and cokernel arguments.
The diagram and its commutativity equations are
\[
 \begin{array}{ccccccccc}
 0&\to&A&\xrightarrow{i}&B&\xrightarrow{p}&C&\to&0\\
 &&\downarrow a&&\downarrow b&&\downarrow c&&\\
 0&\to&A'&\xrightarrow{i'}&B'&\xrightarrow{p'}&C'&\to&0,
 \end{array}
 \qquad bi=i'a,\quad p'b=cp.
\]
\end{hbtheorem}

\begin{hbtheorem}{thm:p10-ab-snake}{proved=annex:p10.ab.11}{The snake connecting map, construction and exactness}
For the same diagram with no invertibility assumption on its vertical
maps, there is a natural exact sequence
\[
 0\to\operatorname{Ker}a\to\operatorname{Ker}b\to
 \operatorname{Ker}c\xrightarrow{\delta}\operatorname{Cok}a\to
 \operatorname{Cok}b\to\operatorname{Cok}c\to0.
\]
To specify the sign and construction, let \(k:K\to C\) be a kernel
of \(c\), form \(P=B\times_C K\) with \(u:P\to B\), \(v:P\to K\),
and let \(t:P\to A'\) be determined by \(i't=bu\).
For \(q:A'\to\operatorname{Cok}a\), the unique connecting map satisfies
\(\delta v=qt\). The proof establishes existence, uniqueness, all
exactness positions, and naturality by the preceding generalized-arrow
method. This is a separately reconstructed diagram result; no module
hypothesis or source embedding theorem is used to justify it.
\end{hbtheorem}

\begin{hbtheorem}{thm:p10-ab-nine}{proved=annex:p10.ab.12}{The nine lemma}
Consider a commutative three-by-three diagram with three short exact
columns and a short exact middle row. The top row is short exact iff
the bottom row is short exact. This is Freyd's statement 2.65, p.~58;
the proof here uses the established snake construction and its categorical
dual, rather than leaving the outer rows to informal element chasing.
\end{hbtheorem}

\begin{hbtheorem}{thm:p10-ab-noether}{proved=annex:p10.ab.13}{Subquotients and the Noether isomorphisms}
For subobjects \(A\subseteq B\subseteq C\),
\[
 0\to B/A\to C/A\to C/B\to0
\]
is exact, so \((C/A)/(B/A)\cong C/B\).
For subobjects \(U,V\subseteq X\), let \(U\cap V\) be their
pullback and \(U+V\) the image of \(U\oplus V\to X\). Then
\(U/(U\cap V)\cong(U+V)/V\). These are categorical subobject
quotients, not quotients of unspecified elements; they expand the source's
statements 2.66--2.67, pp.~58--59, whose full proofs are left as exercises.
\end{hbtheorem}

\begin{hbdefinition}{def:p10-ab-exact-functor}{Additive, left exact, right exact, and exact functors}
A functor is additive when each induced hom-set map is a group
homomorphism. A zero-preserving functor between abelian categories is
left exact if it preserves every sequence \(0\to A\to B\to C\),
right exact if it preserves every \(A\to B\to C\to0\), and exact
if it preserves finite exact sequences. We distinguish preservation of
a specified kernel/cokernel arrow from an equality of chosen objects.
\end{hbdefinition}

\begin{hbtheorem}{thm:p10-ab-functors}{proved=annex:p10.ab.14}{Exact-functor tests and a representable example}
A functor between abelian categories is additive iff it preserves zero
and binary biproduct systems. A zero-preserving left exact or right exact
functor is additive. For an additive functor, left exactness is equivalent
to preserving kernels, right exactness to preserving cokernels, and
exactness to preserving short exact sequences, equivalently both kernels
and cokernels. The covariant functor \((T,-)\) to abelian groups is
additive and left exact, but need not be right exact.
\end{hbtheorem}

\begin{hbtheorem}{thm:p10-ab-module-cases}{proved=annex:p10.ab.15}{Concrete modules, a split vector sequence, and a nonzero snake}
For any unital ring \(R\), left \(R\)-modules form an abelian category.
Kernels are submodules, cokernels are quotient modules, and the comparison
\(A/\ker f\to f[A]\) is an isomorphism.
For any field \(k\), the sequence
\[
 0\to k\xrightarrow{t\mapsto(t,-t,0)} k^3
 \xrightarrow{(x,y,z)\mapsto(x+y,z)} k^2\to0
\]
is split exact, with section \((u,v)\mapsto(u,0,v)\).
For \(n\geq2\), put the exact sequence
\(0\to\mathbb Z\xrightarrow n\mathbb Z\to\mathbb Z/n\to0\)
above a second copy of itself, with vertical maps \(n,n,0\).
The connecting map \(\mathbb Z/n\to\mathbb Z/n\) is the identity
with the stated sign convention. The integer sequence is nonsplit;
\((\mathbb Z/n,-)\) fails right exactness on it.
\end{hbtheorem}

\begin{hbtheorem}{thm:p10-ab-countercases}{proved=annex:p10.ab.16}{Why additive or pointed does not already mean abelian}
The category of all groups fails normality of monomorphisms: the subgroup
generated by \((12)\) in \(S_3\) is not a kernel of any group homomorphism.
Pointed sets have zero objects but their wedge coproduct is not their
product; an explicit surjection identifying two non-basepoints is not a
cokernel. The additive category of finitely generated free abelian groups
has an internal monic and epic map \(2:\mathbb Z\to\mathbb Z\) that
is not invertible. Its internal cokernel is zero and its image/coimage
comparison is multiplication by two, not an isomorphism. Thus the
normality hypotheses cannot be discarded or hidden in a definition of image.
\end{hbtheorem}

These finite exactness arguments do not establish the embedding theorem,
torsion/localization theory, enough injectives, or a derived-functor long
exact sequence. Those are additional constructions with their own hypotheses.
In particular a proof for modules below does not serve as the proof of a
general abelian-category diagram theorem above.
